\section{Arquitetura do DRM Transformer}
\label{sec:arquitetura}

O DRM Transformer mantém a estrutura decoder-only: embeddings de tokens,
blocos autoregressivos, normalização final e cabeça de linguagem. A diferença
principal está no bloco de atenção. Em vez de
\begin{equation}
  \operatorname{softmax}(QK^\top/\sqrt{d_k})V,
\end{equation}
o modelo calcula uma atenção por proximidade geodésica:
\begin{equation}
  \operatorname{Attn}_{ij}
  = \operatorname{softmax}_j\left(-\frac{\tilde{d}_{G,ij}^2}{\tau}\right),
\end{equation}
onde $\tau$ é uma temperatura aprendida com piso mínimo.

Queries e keys são primeiro rotacionadas por RoPE e depois projetadas para o
manifold:
\begin{equation}
  q_m = \sigma(W_q q), \qquad k_m = \sigma(W_k k).
\end{equation}
A sigmoide mantém as coordenadas em $[0,1]$, o que facilita a regularização
toroidal em dimensão quatro. A \emph{MetricNet} recebe $q_m$ e produz
$U(q_m)$, fator low-rank do tensor métrico. A distância usada pela atenção é
então a soma de uma parte euclidiana e uma parte low-rank:
\begin{equation}
  d_G^2(q_m,k_m)=\|q_m-k_m\|^2+
  \|U(q_m)^\top(q_m-k_m)\|^2.
\end{equation}

O campo gravitacional associa uma massa a cada coordenada do manifold e deforma
$U(x)$ por influência baseada em Random Fourier Features. Em termos
interpretativos, tokens com maior massa alteram localmente a geometria de
atenção. Em termos computacionais, isso atua como uma modulação aprendida do
fator métrico.

O \emph{DimensionalGate} aplica uma máscara suave sobre embeddings, permitindo
que a dimensionalidade efetiva varie por token. Esse componente implementa, em
forma neural, a ideia de que $\dim_D(p)$ pode depender do estado local.
